\section{Related Work}


\paragraph{Parallel reasoning and test-time scaling.} Test-time scaling typically follows two paradigms: sequential refinement \citep{madaan2023selfrefine} or parallel generation of multiple reasoning paths \citep{DBLP:journals/corr/abs-2110-14168, setlur2025scalingtesttimecomputeverification, pan2025learning}. While parallel scaling allows for the exploration of diverse solutions, it necessitates a robust mechanism to verify and select the correct output. Existing approaches often rely on access to ground-truth verification signals during inference. For instance, mathematical reasoning often leverages majority voting based on exact answer matching \citep{wang2023selfconsistency, snell2024scalingllmtesttimecompute}, while code generation methods rely on executable test cases or execution feedback \citep{li2025stesttimescaling, jain2025multiturn}. In contrast, we focus on \textit{self-verification}, where the model must judge the quality of its own parallel generations without access to external feedback or ground-truth oracles.


\paragraph{Self-verification and self-aggregation.} Early work demonstrated that LLMs can verify their own Chain-of-Thought (CoT) reasoning~\citep{weng2023selfverification}, though recent studies indicate that pointwise self-verification suffers from a bias toward accepting incorrect solutions~\citep{lu2025doesverificationpayoff}. While sequential self-refinement~\citep{madaan2023selfrefine, stechly2025on} explores verification, it does not address the parallel reasoning setting. Alternatively, self-aggregation methods combine solutions from the same model~\citep{venkatraman2025recursiveselfaggregationunlocksdeep, madaan2025rethinkingthinkingtokensllms} but may suffer from diversity collapse, leading to the loss of correct solutions. Our work addresses these limitations by adopting pairwise self-verification, which mitigates pointwise bias and preserves solution diversity better than aggregation-based approaches.


\paragraph{Generative verifiers and Co-training.} Generative reward models, which produce reasoning before scoring, have been shown to outperform discriminative approaches \citep{mahan2024generative, zhang2025generativeverifiersrewardmodeling}, with pairwise ranking often proving more effective than absolute scoring \citep{jiang2023llmblender, toshniwal2025genselect}. In this context, \citet{zhao2025majority} propose AggLM, which explicitly trains an aggregator via RL to synthesize correct answers for improving majority voting. However, these approaches typically rely on \textit{separate} verifier or aggregator models, incurring significant overhead in compute, memory, and data curation. Recent co-training methods attempt to unify generation and verification to mitigate this cost \citep{sareen2025puttingvaluerlbetter, liu2025trustverifyselfverificationapproach, wang2025coevolvingllmcoderunit}, but they predominantly rely on pointwise verification rewards. Our \method{}-PairRL framework advances this by co-training a single model for \textit{pairwise} self-verification, eliminating the need for external verifiers while enabling efficient online learning.

See Appendix \S\ref{extended_related_works} for more discussion about prior work and extension to this section.

